% =====================================================================
% tikz-closure-flexibility-map-v2.tex
% ---------------------------------------------------------------------
% v2 of the closure flexibility map. Changes vs. v1:
%   * ~20% narrower: every label takes a symmetric two-line break, so the
%     constant-size boxes shrink (2.7cm wide, pitch 3.2cm; v1: 3.4 / 4.0).
%   * a new double arrow between "fixed real wage" and "Leontief closure"
%     (the equivalence of Section 3: the fixed real wage IS the Leontief
%     result once production is rigid), replacing the plain link there.
% Grid unchanged: row 1 = aggregate, row 2 = sectoral; columns = increasing
% flexibility. The elastic-labour-supply box still spans columns 2-3.
% Requires (already in the preamble): \usepackage{tikz} and
%   \usetikzlibrary{arrows.meta, positioning, calc}
% =====================================================================
\begin{figure}[htbp]
\centering
\begin{tikzpicture}[
    node distance = 1.4cm and 2.0cm,
    box/.style = {
      draw, rectangle, rounded corners=3pt,
      minimum width=2.7cm, minimum height=1.1cm,
      align=center, font=\sffamily\small,
      fill=gray!20, line width=0.5pt
    },
    widebox/.style = {
      box, minimum width=5.9cm          % spans two columns + one gap
    },
    placeholder/.style = {
      draw, rectangle, rounded corners=3pt,
      minimum width=2.7cm, minimum height=1.1cm,
      align=center, font=\sffamily\small,
      draw=gray!50, densely dashed, line width=0.4pt
    },
    rowlabel/.style = {font=\sffamily\small},
    axlabel/.style  = {font=\sffamily\small},
    link/.style     = {draw, thick},
    arr/.style      = {-{Stealth}, thick},
    flexarrow/.style = {{Stealth}-{Stealth}, thick}
  ]

  % ---- geometry -------------------------------------------------------
  \def\dx{3.2}    % horizontal pitch between column centres
  \def\rowy{1.9}  % vertical distance between row centres

  % ---- cells: row 1 = aggregate --------------------------------------
  \node[box] (r1c2) at (\dx,0)    {conventional\\CGE};
  \node[box] (r1c3) at (2*\dx,0)  {fixed real\\wage};
  \node[box] (r1c4) at (3*\dx,0)  {Leontief\\closure};

  % ---- cells: row 2 = sectoral ---------------------------------------
  \node[box]         (r2c1) at (0,-\rowy) {no\\reallocation};
  \node[widebox]     (ELS)  at (1.5*\dx,-\rowy) {elastic labor\\supply models};
  \node[placeholder] (r2c4) at (3*\dx,-\rowy) {};

  % ---- connector line through the left aggregate pair (left-right) ----
  \draw[link] (r1c2.east) -- (r1c3.west);

  % ---- double arrow: fixed real wage <-> Leontief closure --------------
  \draw[flexarrow] (r1c3.east) -- (r1c4.west);

  % ---- double arrow: no reallocation -> conventional CGE
  %      up first, then right (one 90-degree corner)
  \draw[flexarrow] (r2c1.north) -- (0,0) -- (r1c2.west);

  % ---- short double arrows of the elastic-labor-supply box -----------
  \draw[flexarrow] (ELS.west)     -- (r2c1.east);        % to no reallocation
  \draw[flexarrow] (2*\dx,-1.35)  -- (r1c3.south);       % up to fixed real wage

  % ---- double arrow: single-wage elastic supply <-> conventional CGE --
  %      (the missing link: with ONE common wage the supply equation
  %      coincides with full employment for every elasticity)
  \draw[flexarrow] (2.6,-1.35) -- (2.6,-0.55);

  % ---- outer axes -----------------------------------------------------
  % y-axis with the two row labels (rotated to avoid overlap)
  \draw[arr] (-2.4,-2.7) -- (-2.4,1.0);
  \draw (-2.4,0)        -- (-2.22,0);
  \draw (-2.4,-\rowy)   -- (-2.22,-\rowy);
  \node[rowlabel, rotate=90] at (-2.75,0)      {aggregate};
  \node[rowlabel, rotate=90] at (-2.75,-\rowy) {sectoral};

  % x-axis: flexibility, increasing to the right
  \draw[arr] (-2.4,-3.2) -- (11.55,-3.2);
  \node[axlabel, below] at (4.58,-3.25) {flexibility};

\end{tikzpicture}
\caption{A map of closure choices along the short-run (labor-market)
  axis. Columns order the closures by increasing flexibility; the
  aggregate row collects the single-wage closures, the sectoral row the
  sector-level ones. The double-width box holds the elastic-labor-supply
  family, tied by short double arrows to the no-reallocation and
  fixed-real-wage corners it interpolates between, and in its single-wage
  form coinciding with the conventional CGE closure (the left-hand vertical
  double arrow); the right-angled double
  arrow marks the reallocation margin from no reallocation up to the
  conventional CGE closure, and the arrow between the fixed real wage and
  the Leontief corner marks their equivalence under demand-only shocks,
  at any substitution elasticities.}
\label{fig:flexibility-map}
\end{figure}